\subsubsection{Operators}
\begin{itemize}
	\item \bi{Union} $\cup$: $x \in R_1 \cup R_2 \Leftrightarrow x \in R_1 \lor x \in R_2$
	      (All tuples from both sets (Only valid for same schemas))
	\item \bi{Difference} $-$: $x \in R_1 - R_2 \Leftrightarrow x \in R_1 \land x \notin R_2$
	      (Tuples that appear in $R_1$, but not in $R_2$)
	\item \bi{Intersection} $\cap$: $x \in R_1 \cap R_2 = R_1 - (R_1 - R_2)$
	      (Tuples that don't appear in both sets)
	\item \bi{Selection} $\sigma$: $x \in \sigma_c(R) \Leftrightarrow x \in R \land c(x) = \texttt{true}$, with $c$ a predicate on the passed in set.
	      (Tuples that fulfil the predicate $c$)
	\item \bi{Projection} $\Pi$: $\Pi_{A_1, \ldots, A_n}(R)$
	      (Keep only a subset of the columns, e.g. for columns \texttt{name}, \texttt{pid}, $\Pi_{\texttt{name}}(R)$ returns only the column \texttt{name})
	\item \bi{Cartesian Product} $\times$: $(x, y) \in R_1 \times R_2 \Leftrightarrow x \in R_1 \land y \in R_2$
	      (Primarily used to express join operations. This however simply joins together the two tables (i.e. similar to expanding terms in maths))
	\item \bi{Renaming} $\rho$: $\rho_{B_1, \ldots, B_n}(R)$
	      (Change the name of the attributes of $R$ to $B_i$)
	\item \bi{Natural Join} $\bowtie$: $R_1(A, B) \bowtie R_2(B, C) = \Pi_{A, B, C} (\sigma_{R_1.B = R_2.B}(R_1 \times R_2))$.
	      (Join two relations into a table on a column. Natural join joins on columns with same name in both (or all) relations)
	      Edge cases:
	      \begin{itemize}
		      \item No shared attributes $R(A, B, C)$, $S(D, E)$: $R \bowtie S = R \times S$
		      \item All attributes shared $R(A, B, C)$, $S(A, B, C)$: $R \bowtie S = R \cap S$
	      \end{itemize}
	\item \bi{Theta Join} $\bowtie_\theta$: $R_1 \bowtie_\theta R_2 = \sigma_\theta(R_1 \times R_2)$
	      (Join with custom predicate $\theta$)
	\item \bi{Equi-Join} $\bowtie_{A = B}$: $R_1 \bowtie_{A = B} = \sigma_{A = B}(R_1 \times R_2)$
	      (Join column A with column B)
	\item \bi{Semi-Join} $\ltimes_C$: $R_1 \ltimes_C R_2 = \Pi_{A_1, \ldots, A_n}(R_1 \bowtie_C$, with $R_1(A_1, \ldots, A_n)$ and $R_2(B_1, \ldots, B_m)$
	      (Returns columns only from one side if there is a match in the join)
	\item \bi{Relational division} $\div$: $R \div S = \Pi_{R - S} R - \Pi_{R - S}((\Pi_{R - S} R) \times S - R)$.
	      In other words, $R \div S = T$, with $T$ being the \textit{largest} relation such that $S \times T \subseteq R$.
	      (Find all rows that do not fulfil a condition)
\end{itemize}
% TODO: Add other join operations

\shade{purple}{Semi-Join Reduction} A useful trick with semi-joins. It is especially useful in distributed DB, where $R$ and $S$ are on different machines.
Suppose we want to Join $R(A, B)$ with $S(B, C)$ on $B$, then we can do the following $R \bowtie S = (R \ltimes \Pi_B S) \bowtie S$.

The idea is to send the column on which the relations are to be joined from the second system to the first,
there execute a semi-join (thus only returning the contents of $R$ that matched with contents of $S$),
then sending only that data to the second machine to finish the full join operation.

Of course, in the real world, users of DB systems don't write relational algebra, as in this case, the order of operations matters for performance,
whereas we want database systems to figure this out by themselves.
